\section{Introdução e delimitação do trabalho}
\label{sec:introducao}

O acompanhamento das condições de armazenamento de grãos pode apoiar a
investigação de alterações no ambiente de conservação. Temperatura e umidade
estão entre as variáveis relevantes nesse contexto, mas uma medição ambiental
isolada não identifica o agente causador de uma alteração. A Organização
Mundial da Saúde descreve a relação entre determinadas condições de
armazenamento, crescimento de fungos e ocorrência de micotoxinas, sem
equiparar essas três situações \cite{who2023}. No presente trabalho, a
investigação é restrita à arquitetura de um protótipo de monitoramento e à
execução local de um classificador demonstrativo.

O MS-Monitor integra uma Heltec WiFi LoRa 32 V2, um DHT22, um termopar tipo K
com conversor MAX6675 e um módulo de vibração SW-420. O firmware adquire as
medições, prepara quatro atributos térmicos e higrométricos e executa um
Random Forest na ESP32. As medições e a classificação são transmitidas a
uma API local para registro em CSV e apresentação em uma interface web.
O sensor de vibração produz telemetria independente; seus dados não são
entradas do modelo atual.

\subsection{Relação com a atividade do Grupo 03}

A especificação original do Grupo G3 solicita a representação dos módulos
de aquisição, processamento, IA, comunicação e visualização, além de um
fluxograma, características técnicas, exemplo de transmissão e plano de
testes \cite{freitas2026}. Naquele enunciado, a IA está associada ao estado
operacional da ventilação. A proposta aqui documentada foi redirecionada
pela equipe ao monitoramento das condições de armazenamento.

Este relatório utiliza os entregáveis arquiteturais da atividade como
referência e explicita essa mudança de objeto. Não apresenta um classificador
de ventilação, nem inclui motor, relé ou atuação automática como partes
implementadas. Consequentemente, não atribui ao SW-420 a capacidade de
determinar, por uma contagem isolada, se um ventilador está ligado.

\subsection{Objetivo e limites de interpretação}

O objetivo é descrever e avaliar a coerência técnica do caminho dos dados,
desde os sensores até a visualização, distinguindo implementação em software,
verificação computacional e validação experimental. Especificamente, o
trabalho apresenta as ligações e a alimentação, o pré-processamento
efetivamente implementado, a exportação do modelo para C++, o contrato da API
e um plano de testes físicos.

A hipótese de pesquisa é investigar se medições térmicas e de umidade podem
contribuir para a classificação de condições associadas à deterioração.
\textbf{O protótipo não demonstrou detecção de fungos, identificação de
espécies, presença de micotoxinas ou segurança sanitária dos grãos.}
Temperatura na ponta da sonda não caracteriza toda a massa armazenada;
umidade relativa do ar também não equivale à umidade do grão ou à atividade
de água. A confirmação dessas condições exige referências e procedimentos
experimentais adicionais.

As descrições técnicas foram confrontadas com o código do projeto, a
documentação do modelo e os resultados computacionais disponíveis na revisão
\texttt{07b1400} \cite{msmonitor2026}. Não se pressupõe que a existência do
firmware compilável comprove, por si só, calibração ou desempenho em bancada.
